\documentclass[10pt,a4paper]{amsart}
\usepackage[margin=19mm]{geometry}
\usepackage{amsmath,amssymb,mathtools}
\usepackage[scr=rsfs,cal=euler]{mathalfa}
\usepackage{xfrac}
\usepackage[unicode,hidelinks,bookmarks=false]{hyperref}
\newcommand{\ie}{i.e.\ }
\newcommand{\cf}{cf.\ }
\newcommand{\resp}{respectively\ }
\newcommand{\Subsection}{\subsection}
\providecommand{\nobreakdash}{\nobreak\hbox{-}}
\setlength{\emergencystretch}{2em}
\multlinegap0pt
\usepackage{siunitx}
\usepackage{siunitx}
\usepackage{siunitx}
\usepackage{siunitx}
\usepackage[normalem]{ulem}
\usepackage{bm}
\usepackage{siunitx}
\usepackage{xcolor}
\usepackage{tikz}
\newtheorem{lemma}{Lemma}
\newcommand{\R}{\mathbb R}
\def\cS{{\mathcal{S}}}
\title{On The Linear Convergence of Bregman Proximal Gradient Methods with Applications to Kullback--Leibler regression}
\author{Jonathan Chirinos-Rodr\'iguez, Christian Daniele, C\'edric F\'evotte, Emmanuel Soubies}
\date{Source: arXiv:2607.05539v1 (CC BY 4.0)}
\begin{document}
\maketitle

\noindent\emph{Source and licence.} On The Linear Convergence of Bregman Proximal Gradient Methods with Applications to Kullback--Leibler regression. Version arXiv:2607.05539v1 (CC BY 4.0), distributed by arXiv under the Creative Commons Attribution 4.0 International licence, http://creativecommons.org/licenses/by/4.0/, as recorded in the arXiv licence metadata for this exact version and shown on the abstract page https://arxiv.org/abs/2607.05539v1. Full licence notice: \texttt{licenses/arxiv-2607.05539-CC-BY-4.0.txt}.\par\smallskip

\noindent\textit{Editorial scope.} Counted unit: the article's lemma alpha(h)>0 (source0.tex lines 1458--1463). The article's complete original proof of it is delivered with it (source0.tex lines 1464--1522, one \texttt{proof} environment of 59 lines). No mathematics is written, repaired or re-derived: the statement and the proof are the article's own lines, in the article's own notation, and no retained line is substituted or re-flowed.  No further source-local context is needed: every cross-reference of every retained line resolves inside the two delivered spans.  Nothing was omitted from any retained span, so the package carries no omission record.  No retained line contains a citation, so the wrapper carries no bibliography and no bibliography entry.  No retained line contains a figure environment, an image-inclusion command or drawing code, so no image is delivered and the package declares no asset dependency.  Editorial formatting only, in addition: an AMS \texttt{amsart} preamble that restates the article's own declarations and macros used by the retained lines, namely the environments \texttt{lemma} on separate counters, together with the standard packages those lines need.  The retained environments print in this document as Lemma 1 in order of appearance, whereas the article prints the same items with the numbering of its own full document.  The complete native source of the article as distributed in the arXiv e-print for this exact version is bundled unchanged as \texttt{sources/204548/source0.tex}.
\begin{lemma} \label{thm: alpha(h)>0}
Let $\phi_\xi:=- \log(x+\xi)$, $\xi>0$, be the one-dimensional smoothed Burg's entropy. Then, for any compact set $\mathcal{S}\subset \mathbb{R}_{\geq 0}$, it holds that
\begin{equation}
\alpha_{\mathcal{S}}(\phi_\xi):=\inf \left\{\frac{d_{\phi_\xi} (x,y)}{d_{\phi_\xi}(y,x)} \mid x\in \mathbb{R}_{\geq 0}, \ y \in \mathcal{S}, x\neq y\right\} >0.
\end{equation}
\end{lemma}

\begin{proof}
Note that the definition of $\alpha_{\mathcal{S}}(\varphi)$ can be rewritten as
$$
\alpha_{\mathcal{S}}(\varphi)=\inf_{x\in\R_{\geq 0}} m(x), \quad \text{where } m(x)=\inf_{y\in \cS} \Delta_x(y),
$$
where, for all $x\in\R_{\geq 0}$,
$$
\Delta_x\colon \R\to \R_{\geq 0}; \quad \Delta_x(y):=
\begin{cases}
\frac{d_{\phi_\xi}(x, y)}{d_{\phi_\xi}(y, x)} & \text{ if } x\neq y\\
1, & \text{ if } x=y.
\end{cases}
$$
To prove the result, we start by showing that $m$ is well defined and continuous. To do so, we first claim that, for all $x\in \R_{\geq 0}$, $\Delta_x$ is continuous. First, if $x\neq y$, then it is continuous as the division of two Bregman divergences, the denominator being nonzero everywhere. It is therefore sufficient to prove that 
$$
\lim_{x\to y}\frac{d_{\phi_\xi}(x, y)}{d_{\phi_\xi}(y, x)}=1.
$$
We have that
$$
\lim_{x\to y}\frac{d_{\phi_\xi}(x, y)}{d_{\phi_\xi}(y, x)}=\lim_{x\to y}\frac{\log(\frac{y+\xi}{x+\xi})+\frac{x-y}{y+\xi}}{ \log(\frac{x+\xi}{y+\xi})+\frac{y-x}{x+\xi}}=\frac{\log(t)+1/t-1}{-\log(t)+t-1},
$$
where simply $t:=(y+\xi)/(x+\xi)$. In this case, as $x\to y$ is equivalent to $t\to 1$, we need to prove that
$$
\lim_{t\to 1}\frac{\log(t)+1/t-1}{-\log(t)+t-1}=1.
$$
By applying L'H\^opital's rule twice, we derive that
$$
\begin{aligned}
\lim_{t\to 1} \frac{\log(t)+1/t-1}{-\log(t)+t-1}=\lim_{t\to 1}\frac{t\log(t)+1-t}{-t\log(t)+t^2-t}= \lim_{t\to 1} \frac{\log(t)}{-\log(t)+2t-2}=\lim_{t\to 1} \frac{1/t}{-1/t+2}=1,
\end{aligned}
$$
concluding the proof of the claim. Now, as $\Delta_x(y)>0$ for all $y\in\cS$, $\cS$ is compact, and $\Delta_x$ is continuous, we get by the extreme value theorem that the minimum over $\mathcal{S}$ is attained, and so $m(x)=\min_{y\in \cS} \Delta_x(y)>0$ for all $x\in \R$. Additionally, as $m$ is defined as the minimum of continuous functions over a compact set, it is continuous too. Thus to conclude it suffices to show that $m$ is coercive. We want to show that 
$$
    \lim_{x\to\infty} \ \underset{y \in \cS}{\inf}\, \frac{d_{\phi_\xi}(x,y)}{d_{\phi_\xi}(y,x)} = \underset{x \to + \infty}{\lim} \ \underset{y \in \cS}{\inf} \, \frac{\log(\frac{y+\xi}{x+\xi})+\frac{x-y}{y+\xi}}{ \log(\frac{x+\xi}{y+\xi})+\frac{y-x}{x+\xi}}=+\infty.
$$
To do so, we will use compactness of $\mathcal{S}$ to lower bound (resp. upper bound) the numerator (resp. denominator) of the ratio above in such a way that it does not depend on $y$ anymore. We start observing that, since $\mathcal{S}$ is compact, there exist $0<\delta_1\leq\delta_2<+\infty$ such that $y\in[\delta_1,\delta_2]$. It follows that $\log(y+\xi)\geq \log(\delta_1)$ and
$$
\frac{x-y}{y+\xi}\geq \frac{x-y}{\delta_2+\xi}\geq\frac{x-\delta_1}{\delta_2+\xi}.
$$
Then, the numerator satisfies,
$$
\begin{aligned}
d_{\phi_\xi}(x,y)\geq  -\log(x+\xi)+\log(\delta_1)+\frac{x-\delta_1}{\delta_2+\xi}.
\end{aligned}
$$
On the other hand, we have that $\frac{y-x}{x+\xi}\leq \frac{y}{x+\xi}\leq \frac{\delta_2}{\xi}$, and so, the denominator satisfies
$$
 d_{\phi_\xi}(y,x)\leq -\log(\xi)+\log(x+\xi)+\frac{\delta_2}{\xi}.
$$
Combining both estimates, we have that
\begin{equation}\label{eq:lower_bound_m}
m(x)=\underset{y \in \cS}{\inf} \, \frac{d_{\phi_\xi}(x,y)}{d_{\phi_\xi}(y,x)}\geq\frac{-\log(x+\xi)+\log(\delta_1)+\frac{x-\delta_1}{\delta_2+\xi}}{-\log(\xi)+\log(x+\xi)+\frac{\delta_2}{\xi}},
\end{equation}
where the right hand side above clearly satisfies that
$$
\underset{x \to + \infty}{\operatorname{lim}} \frac{-\log(x+\xi)+\log(\delta_1)+\frac{x-\delta_1}{\delta_2+\xi}}{-\log(\xi)+\log(x+\xi)+\frac{\delta_2}{\xi}}= + \infty.
$$
Taking limits on both sides of \eqref{eq:lower_bound_m} leads to the desired result. 
\end{proof}

\end{document}
